\problemname{Hasty Haul}

\illustration{0.35}{tables-small.jpg}{
    The registration area at NWERC 2024.
    Photo by Maarten Sijm
}

Two teams participating in NWERC have hatched a devious plan:
% to mess with the organizers.
they want to get in the organizers' heads and mess with them a bit, but just subtly enough
that they will not get caught.

The plan is as follows: during the contest registration, where they get
their shirts and goodies, they will first wait for the organizer to lean down to grab
a shirt.
Then, they will quickly lift up one of the $k$ pieces of furniture and move it around.
With the limited time
frame and the need to do it silently, each team can only manage to move one piece of furniture.

To make sure that they do not get caught, each piece of furniture has to be back at its original place at the end of the registration.
From previous visits, the teams know the dimensions of the room and the amount of furniture in it
ahead of time, but they do not know the current arrangement of the furniture.
To keep things simple, the two teams want to use the same strategy.
This means that both teams will make the same move when they encounter the same layout.
After deciding on a common strategy and setting off for the contest, they will not be able to communicate.
Thus, they plan to come up with a strategy that ensures that, regardless of the
arrangement of the furniture,
the team that arrives later will undo whatever the team that arrived first did.
For some combinations of room size and amount of furniture in the room this is impossible,
so pulling off this kind of stunt would be risky.

At least they know for sure they are the only ones pulling this kind of stunt, so no furniture
should move around between the two teams' arrivals. Help them find a strategy for their prank!

\begin{Input}
    The input consists of:
    \begin{itemize}
        \item One line with an integer $t$ ($1\leq t\leq 10\,000$), the number of test cases.
        \item For each test case, the input consists of:
        \begin{itemize}
        	\item One line with three integers $h$, $w$, and $k$ ($1\leq h,w\leq 8$, $1 \leq k < h \cdot w$),
                the height and width of the registration room and the number of pieces of furniture in the room.
            \item $h$ lines with $w$ characters, each character being either `\texttt{.}' or `\texttt{\#}', the state of the room.
                A `\texttt{.}' represents an empty area, and `\texttt{\#}'
                represents a single movable piece of furniture.
                It is guaranteed that the room
                contains exactly $k$ pieces of furniture.
        \end{itemize}
        It is guaranteed that there is at least one piece of furniture and some empty area.
    \end{itemize}
    This is a multi-pass problem. Your program will be invoked multiple times,
    possibly more than twice.
    Your program must act consistently within each invocation, but also across
    invocations.

    For testing purposes, the number and input of subsequent passes will depend
    on the output of your submission.

    A testing tool is provided to help you develop your solution.
\end{Input}

\begin{Output}
    For each test case, output ``\texttt{risky}'' if there is no valid strategy the teams could come up with
    beforehand that works regardless of how the $k$ pieces of
    furniture end up being placed.
    %% This means there is no plan that makes sure all the requirements above are
    %% satisfied, no matter how the furniture in this test case is arranged at the start.
    Otherwise, specify how a piece of furniture is moved:
    first output the position of the piece of furniture you move and then output the position where it should be moved to.
    Both positions must first specify the row $r$ ($1\leq r \leq h$, counting
    from the top) and then
    the column $c$ ($1\leq c\leq w$, counting from the left).
\end{Output}
